% !TEX root = ../../main.tex

\section{The Black-Scholes-Merton Model}
The Black-Scholes-Merton model provides a theoretical estimate of the price of European-style options. It assumes that the price of the underlying asset follows a Geometric Brownian Motion (GBM).

\subsection{Assumptions}
\begin{itemize}
    \item \textbf{Efficient Markets:} Asset prices follow a log-normal distribution.
    \item \textbf{No Dividends:} The underlying stock pays no dividends during the option's life.
    \item \textbf{Constant Parameters:} Risk-free rate ($r$) and volatility ($\sigma$) are constant and known.
    \item \textbf{Frictionless Markets:} No transaction costs or taxes; short-selling is permitted.
    \item \textbf{Continuity:} Trading is continuous, and the price path is continuous (no jumps).
\end{itemize}

\subsection{The Black-Scholes PDE}
The value of any derivative $V(S,t)$ must satisfy the following second-order partial differential equation:
\begin{equation}
    \frac{\partial V}{\partial t} + rS \frac{\partial V}{\partial S} + \frac{1}{2}\sigma^2 S^2 \frac{\partial^2 V}{\partial S^2} - rV = 0
\end{equation}
\textbf{Intuition:} The PDE is derived through a "delta-hedging" argument. By creating a portfolio of the option and a specific amount of the underlying stock, one can eliminate risk (stochasticity) locally, meaning the portfolio must earn the risk-free rate.

\subsection{Closed-Form Solution (European Options)}
For a European call ($C$) and put ($P$), the prices are:
\begin{align}
    C(S, t) &= S_t N(d_1) - Ke^{-r(T-t)} N(d_2) \\
    P(S, t) &= Ke^{-r(T-t)} N(-d_2) - S_t N(-d_1)
\end{align}
Where $N(\cdot)$ is the cumulative distribution function of the standard normal distribution, and:
\begin{align}
    d_1 &= \frac{\ln(S_t/K) + (r + \frac{\sigma^2}{2})(T-t)}{\sigma\sqrt{T-t}} \\
    d_2 &= d_1 - \sigma\sqrt{T-t}
\end{align}

\subsection{Understanding $N(d_1)$ and $N(d_2)$}
The call formula can be seen as $C = [\text{Expected Value of Stock Receivied}] - [\text{Expected Cost of Exercise}]$.
\begin{itemize}
    \item \textbf{$N(d_2)$ (The Probability):} This is the risk-neutral probability that the option will expire in-the-money ($S_T > K$). If $N(d_2) = 0.5$, there is a 50\% chance you will exercise the option.
    \item \textbf{$N(d_1)$ (The Replicator):} This is the \textit{Delta} ($\Delta$). It represents the factor by which the stock's price change affects the option's value. 
\end{itemize}
\textbf{Intuition Example:} Imagine an option very deep in-the-money. $N(d_2)$ is nearly 1 (certainty of exercise). $N(d_1)$ is also nearly 1 (the option moves exactly like the stock). Conversely, for an out-of-the-money option, $N(d_2)$ might be 0.1, but $N(d_1)$ might be 0.15. The Delta is always slightly higher than the probability of exercise because it accounts for the "magnitude" of the potential gain.

\subsection{Dynamic Hedging}
The BSM model assumes you can create a risk-free portfolio by constantly rebalancing.
\begin{itemize}
    \item \textbf{Delta Neutrality:} If you sell 1 Call with $\Delta=0.6$, you must buy 0.6 shares of the stock. Your net Delta is $0$.
    \item \textbf{The Cost of Hedging:} Because Gamma exists, as the stock goes up, your Delta increases. To stay hedged, you must buy more stock at a higher price. If the stock goes down, you must sell stock at a lower price. This "buy high, sell low" behavior is the reason why being "Short Gamma" (selling options) is risky and requires a premium (the option price).
\end{itemize}

\subsection{The Greeks}
The Greeks measure the sensitivity of the option price to various model inputs.

\begin{table}[h!]
\centering
\begin{tabular}{|l|l|l|}
\hline
\textbf{Greek} & \textbf{Definition} & \textbf{Formula (Call)} \\ \hline
Delta ($\Delta$) & $\frac{\partial V}{\partial S}$ & $N(d_1)$ \\ \hline
Gamma ($\Gamma$) & $\frac{\partial^2 V}{\partial S^2}$ & $\frac{N'(d_1)}{S\sigma\sqrt{T-t}}$ \\ \hline
Vega ($\nu$) & $\frac{\partial V}{\partial \sigma}$ & $S\sqrt{T-t} N'(d_1)$ \\ \hline
Theta ($\Theta$) & $-\frac{\partial V}{\partial \tau}$ & $-\frac{S N'(d_1) \sigma}{2\sqrt{T-t}} - rKe^{-rt}N(d_2)$ \\ \hline
Rho ($\rho$) & $\frac{\partial V}{\partial r}$ & $K(T-t)e^{-r(T-t)}N(d_2)$ \\ \hline
\end{tabular}
\caption{Summary of First and Second Order Greeks}
\end{table}

\subsubsection{Deep Dive: Delta and Gamma}
\begin{itemize}
    \item \textbf{Delta ($\Delta$):} For a call option, $\Delta \in [0, 1]$. It represents the probability of the option expiring in-the-money (ITM) in a risk-neutral world. It is also the number of shares needed to hedge one option.
    \item \textbf{Gamma ($\Gamma$):} Measures the stability of Delta. High Gamma (usually near-the-money and close to expiry) indicates that Delta changes rapidly, making the portfolio difficult and expensive to rebalance.
\end{itemize}

\subsection{Put-Call Parity}
A fundamental relationship between the price of a European call and a European put with the same strike and expiry:
\begin{equation}
    C - P = S_t - Ke^{-r(T-t)}
\end{equation}
If this equality does not hold, an arbitrage opportunity exists.

\subsection{Implied Volatility and the Smile}
The BSM model assumes constant volatility. However, in practice, if we take the market price of an option and "back out" the volatility using the BSM formula, we find that $\sigma$ varies by strike price and maturity. This phenomenon is known as the \textbf{Volatility Smile} or \textbf{Skew}.